\scanpage{102}

\mspage{108}

Mais cela reste-il vrai après tout changement de base ?

La situation semble ici meilleure pour $f$ lisse (au lieu de propre), \deleted{alors} \inserted{en interprétant} les conditions a) et b) sous forme homologique, comme signifiant que $f_!(e_X)$ est un objet \uline{loc constant} de $\mathrm{HOT}(Y)$, qui se calcule fibre par fibre.

Comme $f$ lisse implique que $f_!$ commute à tout changement de base, \deleted{cette} cette hyp. \deleted{que $f$ \unsure{soit}} sur $f$ (que $f$ soit une HOT-fibration \deleted{\unsure{co}}homologique, peut-être pas complète) \deleted{implique} est stable par chgt de base, donc implique que c'est une HOT fibration cohomologique complète.

Mais cet argument utilise \deleted{des \unsure{opérations}} des propriétés générales des dérivateurs pour HOT, que je n'ai pas encore établies -- surtout pas dans le domaine $\mathrm{Cat}$ tout entier ! Donc il est sain de se méfier. Je préfère tâter la difficulté, en restant d'abord dans le cadre des \unsure{questions|prénotions|propriétés|opérations} \uline{cohomologiques} (\unsure{au moins|au mouv.}), dont j'ai une certaine maîtrise. Je suppose un $f$ propre encore, et \deleted{les conditions}\inserted{que les}

\[
X_y\to X_{/y},\quad X_{/y}\to X_{/\unsure{y|y'}},
\]

% archive sheet number: 102
